\documentclass{article}

\usepackage{amsmath}
\usepackage{listings}
\usepackage{amsopn}

\lstdefinelanguage{Rust}{
    morekeywords={fn, let, mut, pub, struct, use, mod, if, else, for, return, Some, None},
    sensitive=true,
    morecomment=[l]{//},
    morestring=[b]"
}
\usepackage{amssymb}



\begin{document}

    \section{Theoretische Grundlagen}
    \label{sec:grundlagen}

    Dieses Kapitel legt die mathematischen und physikalischen Grundlagen dar,
    auf denen die Simulationspipeline aufbaut.
    Das Verständnis dieser Konzepte ist Voraussetzung für die Bewertung der
    in späteren Kapiteln beschriebenen Implementierungsentscheidungen.

    \subsection{Akustische Grundlagen}
    \label{subsec:akustik}

    \paragraph{Das Inverse Quadratgesetz}

    Im freien, reflexionsfreien Schallfeld breitet sich Schall kugelförmig vom Emitter aus.
    Da die abgestrahlte Schallleistung $P$ über eine mit $r^2$ wachsende Kugeloberfläche
    $A = 4\pi r^2$ verteilt wird, nimmt die Schallintensität $I$ mit dem Quadrat der Entfernung ab:
    \[I(r) = \frac{P}{4\pi r^2}\]
    Da die Schallintensität proportional zum Quadrat des Schalldrucks ist ($I \propto p^2$),
    folgt für den Schalldruck:
    \[p(r) \propto \frac{1}{r}\]
    Eine Verdopplung des Abstands zur Quelle halbiert den Schalldruck, was einer
    Pegelminderung von $6\,\text{dB}$ entspricht.
    In der Simulation wird dieser geometrische Energieverlust implizit berücksichtigt:
    Strahlen, die einen längeren Weg zur Mikrofonkapsel zurücklegen, treffen mit einem
    entsprechend geringeren Druckbeitrag ein.

    \paragraph{Materialabsorption}

    Trifft ein Schallstrahl auf eine Raumgrenze, wird ein Teil der Schallenergie absorbiert.
    Dieser Vorgang wird durch den dimensionslosen Absorptionskoeffizienten $\alpha \in [0{,}1]$
    beschrieben. Der Restdruck nach einer einzelnen Reflexion berechnet sich zu:
    \[p_{\text{ref}} = p_{\text{ein}} \cdot (1 - \alpha)\]
    Bei einer Sequenz von $n$ Reflexionen mit den jeweiligen Koeffizienten
    $\alpha_1, \ldots, \alpha_n$ gilt:
    \[p_n = p_0 \cdot \prod_{k=1}^{n}(1 - \alpha_k)\]
    Ein vollständig absorbierendes Material ($\alpha = 1{,}0$) vernichtet den Strahl bei der
    ersten Reflexion, während ein perfekt reflektierendes Material ($\alpha = 0{,}0$) keine
    Energie entzieht. Die Simulation terminiert Strahlen, sobald ihr verbleibender Druckwert
    den konfigurierbaren Grenzwert \texttt{min\_pressure} unterschreitet.

    \paragraph{Frühe Reflexionen und späte Nachhallzeit}

    Die akustische Impulsantwort eines Raumes lässt sich zeitlich in zwei charakteristische
    Phasen unterteilen.
    Innerhalb der ersten ca. $80\,\text{ms}$ nach dem Direktschall treffen die
    \emph{frühen Reflexionen} (Early Reflections) ein: diskrete, noch einzeln wahrnehmbare
    Echos über klar definierte, kurze Ausbreitungspfade, die das Raumgefühl und die
    Tiefenwahrnehmung maßgeblich prägen.
    Mit zunehmender Zeit überlagern sich Reflexionen höherer Ordnung zu einem dichten,
    statistisch gleichverteilten Nachhallfeld. Diese \emph{späte Nachhallzeit}
    (Late Reverberation) klingt exponentiell ab und wird durch die Nachhallzeit $RT_{60}$
    charakterisiert -- der Zeitraum, in dem der Schalldruckpegel um $60\,\text{dB}$ abfällt.
    Das stochastische Raytracing erfasst beide Phasen natürlich, da sowohl kurze Direktpfade
    als auch komplexe Mehrfachreflexionen in der Impulsantwort akkumuliert werden.

    \subsection{Monte-Carlo-Integration und Raytracing}
    \label{subsec:montecarlo}

    \paragraph{Das akustische Integrationsproblem}

    Der Schalldruck am Empfänger ist das Integral über alle möglichen Schallpfade vom Emitter,
    gewichtet nach Weglänge und materialabhängigen Verlusten. Da die Anzahl dieser Pfade durch
    Mehrfachreflexionen kombinatorisch explodiert, ist eine analytische Auswertung für
    beliebige Raumgeometrien nicht praktikabel. Deterministisches Ray Tracing kann zwar exakte
    Pfade verfolgen, skaliert jedoch exponentiell mit der Anzahl der Reflexionen.

    \paragraph{Der Monte-Carlo-Ansatz}

    Die Monte-Carlo-Integration umgeht dieses Problem durch stochastische Stichprobennahme.
    Statt alle Pfade zu berechnen, schätzt man den Erwartungswert eines Integrals durch den
    arithmetischen Mittelwert von $N$ zufällig gezogenen Stichproben:
    \[\mathbb{E}[f(X)] \approx \frac{1}{N} \sum_{i=1}^{N} f(x_i)\]
    Nach dem Gesetz der großen Zahlen konvergiert dieser Schätzer für $N \to \infty$ gegen
    den wahren Integralwert. Der statistische Fehler nimmt dabei mit $\frac{1}{\sqrt{N}}$ ab:
    ein viermal so großes Stichprobenvolumen halbiert den Fehler.

    \paragraph{Anwendung: Stochastisches Raytracing}

    In der Simulation entspricht eine Stichprobe einem Schallstrahl, der in einer zufälligen
    Richtung $\theta \in [0, 2\pi)$ vom Emitter ausgesendet wird:
    \[\vec{D} = (\cos\theta,\; \sin\theta), \quad \theta \sim \mathcal{U}(0,\; 2\pi)\]
    Jeder Strahl verfolgt seinen Ausbreitungspfad, reflektiert an Wänden und registriert
    einen Treffer, sobald er die Mikrofonkapsel schneidet. Der akkumulierte Beitrag dieses
    Treffers -- Laufzeit $\tau$ und Restdruck $p$ -- repräsentiert eine Stichprobe des
    akustischen Integrals. Über $N = 10^6$ bis $10^8$ Strahlen konvergiert die Gesamtheit
    der Treffer zur vollständigen Impulsantwort des Raumes.
    Der wesentliche Vorteil gegenüber deterministischen Verfahren liegt in der trivialen
    Parallelisierbarkeit: Da jeder Strahl unabhängig von allen anderen ist, lassen sich die
    $N$ Stichproben direkt auf beliebig viele CPU-Kerne aufteilen
    (siehe Kapitel~\ref{sec:concurrency}).

    \subsection{Digitale Signalverarbeitung: Impulsantwort und schnelle Faltung}
    \label{subsec:dsp}

    \paragraph{Die Impulsantwort (IR)}

    Ein lineares, zeitinvariantes System -- wie ein Raum aus akustischer Perspektive -- ist
    vollständig durch seine \emph{Impulsantwort} $h(t)$ charakterisiert: die Ausgabe des
    Systems bei einem idealen Dirac-Impuls $\delta(t)$ am Eingang.
    Die von der Rust-Engine berechnete Menge an (Laufzeit, Druck)-Paaren
    $\{(\tau_i, p_i)\}_{i=1}^{M}$ stellt eine diskrete Approximation dieser Impulsantwort dar:
    \[h(t) \approx \sum_{i=1}^{M} p_i \cdot \delta(t - \tau_i)\]
    Das Python-Modul \texttt{convolution.py} diskretisiert diese Darstellung in ein
    sample-genaues Array, indem jede Verzögerung $\tau_i$ in den Sample-Index
    $n_i = \lfloor \tau_i \cdot f_s \rceil$ umgerechnet und Druckwerte akkumuliert werden
    (bei $f_s = 44100\,\text{Hz}$).

    \paragraph{Das Faltungsintegral}

    Das Ausgangssignal $y(t)$ -- das hallbehaftete Audio -- ergibt sich aus der linearen
    Faltung des Eingangssignals $x(t)$ mit der Impulsantwort $h(t)$:
    \[y(t) = (x * h)(t) = \int_{-\infty}^{\infty} x(\tau)\, h(t - \tau)\, d\tau\]
    In der diskreten Signalverarbeitung wird dieses Integral zur Summe:
    \[y[n] = \sum_{k=0}^{M-1} h[k] \cdot x[n - k]\]
    Die direkte Auswertung hat eine Rechenkomplexität von $\mathcal{O}(N \cdot M)$, wobei $N$
    die Länge des Audiosignals und $M$ die Länge der Impulsantwort bezeichnet.
    Bei realen Impulsantworten mit $M$ von mehreren hunderttausend Samples ist eine direkte
    Berechnung unpraktikabel.

    \paragraph{Schnelle Faltung via FFT (Overlap-Add)}

    Das \emph{Faltungstheorem} der Fouriertransformation besagt, dass eine Faltung im
    Zeitbereich einem punktweisen Produkt im Frequenzbereich entspricht:
    \[Y(f) = \mathcal{F}\{x\}(f) \cdot \mathcal{F}\{h\}(f)
    \implies y[n] = \mathcal{F}^{-1}\{Y(f)\}\]
    Durch den Einsatz der Fast Fourier Transform (FFT) sinkt die Komplexität auf
    $\mathcal{O}((N+M)\log(N+M))$, was bei großen $M$ einen massiven Speedup bedeutet.

    Da das Audiosignal wesentlich länger sein kann als eine einzelne FFT erlaubt,
    wird das \emph{Overlap-Add}-Verfahren eingesetzt: Das Signal $x[n]$ wird in
    überlappungsfreie Blöcke der Länge $L$ aufgeteilt, jeder Block wird separat mit $h[n]$
    gefaltet, und die resultierenden Ausgabeblöcke -- je $L + M - 1$ Samples lang -- werden
    additiv überlagert. Die Implementierung nutzt hierfür
    \texttt{scipy.signal.oaconvolve}, welche diesen Algorithmus effizient realisiert und
    sowohl mono- als auch mehrkanalige Signale verarbeitet.


    \section{Physikalischer Kern: 2D-Strahlensegmentschnitt}
    \label{sec:2d_ray_intersection}

    Der rechnerische Kern der Physik-Engine basiert auf robuster linearer Algebra, um Randkollisionen zu lösen.
    Eine genaue mathematische Modellierung in diesem Stadium ist für hochauflösende akustische Simulationen entscheidend.
    Um den Rechendurchsatz zu maximieren, was bei der Berechnung von Millionen von Strahlen zwingend erforderlich ist, verzichtet die Engine auf
    rechnerisch aufwändige trigonometrische Funktionen zugunsten reiner Vektormathematik.

    \subsection{Geometrische Definitionen}\label{subsec:geometric-definitions}

    %The two primary geometric entities in the 2D simulation space are the acoustic ray and the physical boundary (representing a wall).
    Die beiden wichtigsten geometrischen Entitäten im 2D-Simulationsraum sind der akustische Strahl und die physikalische Grenze (die eine Wand darstellt).

    \paragraph{Der Strahl}
    %A sound ray is defined by an origin point $O$ and a normalized direction vector $\vec{D}$.
    %Any point $P$ along the ray's trajectory can be calculated as a function of the scalar distance $t$ traveled along that direction:
    %\[P(t) = O + t\vec{D}\]
    %where $t \ge 0$ ensures the ray propagates forward in space from the emitter.
    Ein Schallstrahl wird durch einen Ursprungspunkt $O$ und einen normalisierten Richtungsvektor definiert $\vec{D}$.
    Jeder Punkt $P$ entlang der Strahlenbahn kann als Funktion der skalaren Distanz $t$ entlang dieser Richtung berechnet werden:
    \[P(t) = O + t\vec{D}\]
    wobei $t \ge 0$ sicherstellt, dass der Strahl vom Emitter aus im Raum weitergeht.

    \paragraph{Die Wand (Die Strecke)}
    %A wall is represented as a finite line segment defined by two endpoints, $A$ and $B$.
    %Any point $W$ on this segment is found by interpolating between the endpoints using a scalar parameter $u$.
    %Defining the wall vector as $\vec{V} = B - A$, the parametric equation is given by:
    %\[W(u) = A + u\vec{V}\]
    %For a point to lie strictly within the physical bounds of the wall segment, the condition $0 \le u \le 1$ must be satisfied.

    Eine Wand wird als Liniensegment dargestellt, das durch zwei Endpunkte, $A$ und $B$, definiert ist.
    Jeder Punkt $W$ auf diesem Segment wird durch Interpolation zwischen den Endpunkten mit einem skalaren Parameter $u$ gefunden.
    Definiert man den Wandvektor als $\vec{V} = B - A$, ergibt sich die parametrische Gleichung durch:
    \[W(u) = A + u\vec{V}\]
    Damit ein Punkt strikt innerhalb der physikalischen Grenzen des Wandsegments liegt, muss die Bedingung $0 \le u \le 1$ erfüllt sein.

    \subsection{Schnittpunktberechnung}\label{subsec:intersection-formula}

    %To determine the point of intersection, the ray and wall equations are equated:
    %\[O + t\vec{D} = A + u\vec{V}\]
    %This forms a system of linear equations to solve for $t$ (propagation distance) and $u$ (segment intersection parameter).
    %By defining a vector $\vec{E} = A - O$ that points from the ray's origin to the wall's starting point, the equation rearranges to:
    %\[t\vec{D} - u\vec{V} = \vec{E}\]
    %In two-dimensional space, this system can be efficiently solved using the 2D analog of the cross product (often referred to as the perp-dot product).
    %For two arbitrary 2D vectors $\vec{v_1} = (x_1, y_1)$ and $\vec{v_2} = (x_2, y_2)$, this yields a scalar determinant: $\vec{v_1} \times \vec{v_2} = x_1 y_2 - y_1 x_2$.
    %
    %By defining the intersection determinant as $\det = \vec{D} \times \vec{V}$, the parameters $t$ and $u$ can be isolated and solved algebraically:
    %\begin{gather*}
    %    t = \frac{\vec{E} \times \vec{V}}{\det}\\
    %    u = \frac{\vec{E} \times \vec{D}}{\det}\\
    %\end{gather*}

    Um den Schnittpunkt zu bestimmen, werden die Strahlen- und Wandgleichungen gleichgesetzt:
    \[O + t\vec{D} = A + u\vec{V}\]
    Dies bildet ein System linearer Gleichungen, um $t$ (Ausbreitungsdistanz) und $u$ (Segmentschnittparameter) zu lösen.
    Durch die Definition eines Vektors $\vec{E} = A - O$, der vom Ursprung des Strahls zum Ausgangspunkt der Wand zeigt, ordnet sich die Gleichung um:
    \[t\vec{D} - u\vec{V} = \vec{E}\]
    Im zweidimensionalen Raum kann dieses System effizient mit dem 2D-Analogon des Kreuzprodukts gelöst werden (oft als Perp-Dot-Produkt bezeichnet).
    Für zwei beliebige 2D-Vektoren $\vec{v_1} = (x_1, y_1)$ und $\vec{v_2} = (x_2, y_2)$ ergibt sich eine skalare Determinante:$\vec{v_1} \times \vec{v_2} = x_1 y_2 - y_1 x_2$.

    Durch die Definition der Schnittdeterminante als $\det = \vec{D} \times \vec{V}$ können die Parameter $t$ und $u$ isoliert und algebraisch gelöst werden:
    \begin{gather*}
        t = \frac{\vec{E} \times \vec{V}}{\det}\\
        u = \frac{\vec{E} \times \vec{D}}{\det}\\
    \end{gather*}

    \subsection{Schnittpunkt-Validierung (Trefferbedingungen)}\label{subsec:intersection-validation-(hit-conditions)}

    %A mathematically valid solution to the linear system does not inherently guarantee a physical collision.
    %The following constraints must be verified to confirm a valid acoustic reflection:
    %\begin{itemize}
    %    \item \textbf{Parallelism:} If $\det \approx 0$, the ray and the wall are perfectly parallel and will never intersect.
    %    In the implementation, a small epsilon threshold is used rather than an exact zero check to prevent floating-point inaccuracies.
    %    \item \textbf{Directionality:} If $t < 0$, the mathematical lines intersect, but the collision occurs behind the ray's point of emission.
    %    \item \textbf{Segment Bounds:} If $u < 0$ or $u > 1$, the ray intersects the infinite line containing the wall, but misses the physical line segment (e.g., simulating sound passing through an open doorway).
    %\end{itemize}
    %A physical boundary collision is only registered if $\det \neq 0$, $t \ge 0$, and $0 \le u \le 1$.

    Eine mathematisch gültige Lösung des linearen Systems garantiert nicht automatisch eine physikalische Kollision.
    Die folgenden Einschränkungen müssen überprüft werden, um eine gültige akustische Reflexion zu bestätigen:
    \begin{itemize}
        \item \textbf{Parallelität:} Wenn $\det \approx 0$, sind der Strahl und die Wand perfekt parallel und schneiden sich niemals
        In der Implementierung ist ein kleiner Epsilon-Schwellenwert anstelle einer exakten Nullprüfung verwendet, um Gleitkommaungenauigkeiten zu vermeiden:

        \begin{lstlisting}[language=Rust]
    if det.abs() < 1e-6 {
        return None;
    }
        \end{lstlisting}

        \item \textbf{Richtungsabhängigkeit:} Wenn $t < 0$, schneiden sich die mathematischen Geraden, aber die Kollision erfolgt hinter dem Aussendungspunkt des Strahls:
        \begin{lstlisting}[language=Rust]
    if ray_t < 0.0 {
        return None;
    }
        \end{lstlisting}
        \item \textbf{Segmentgrenzen:} Wenn $u < 0$ oder $u > 1$, schneidet der Strahl zwar die unendliche Gerade, auf der die Wand liegt, verfehlt jedoch die physische Strecke
        \begin{lstlisting}[language=Rust]
    if wall_u < 0.0 || wall_u > 1.0 {
        return None;
    }
        \end{lstlisting}
    \end{itemize}

    \subsection{Vektorreflexion}\label{subsec:reflection-vector-calculation}

    %Upon confirming a valid intersection, the resulting trajectory of the sound wave must be calculated.
    %This requires the surface normal vector $\vec{N}$, which points strictly perpendicular away from the wall.
    %For a given 2D wall vector $\vec{V} = (x, y)$, the normalized perpendicular normal is $\vec{N} = (-y, x)$.

    Nach der Bestätigung eines gültigen Schnittpunkts muss die resultierende Trajektorie der Schallwelle berechnet werden.
    Dies erfordert den Oberflächennormalenvektor $\vec{N}$, der streng senkrecht von der Wand weg zeigt
    Für einen gegebenen 2D-Wandvektor $\vec{V} = (x, y)$ ist die normierte senkrechte Normale $\vec{N} = (-y, x)$.

    %The reflected ray direction $\vec{R}$ is determined by taking the incoming ray direction $\vec{D}$ and subtracting twice the projection of that direction onto the surface normal:
    %\[\vec{R} = \vec{D} - 2(\vec{D} \cdot \vec{N})\vec{N}\]
    %Within the Rust codebase, this vector reflection calculation is optimized using the \texttt{reflect} method provided by the \texttt{glam} crate,
    %which leverages SIMD instructions to perform this operation with minimal overhead.

    Die Richtung des reflektierten Strahls $\vec{R}$ wird bestimmt, indem von der Richtung des einfallenden Strahls $\vec{D}$ das Doppelte der Projektion dieser Richtung
    auf die Oberflächennormale subtrahiert wird.
    \[\vec{R} = \vec{D} - 2(\vec{D} \cdot \vec{N})\vec{N}\]
    Innerhalb der Rust-Codebasis wird diese Berechnung der Vektorreflexion durch die \texttt{reflect} Methode optimiert, die vom \texttt{glam}-Crate bereitgestellt wird und
    Single Instruction, Multiple Data (SIMD)-Instruktionen nutzt, um diese Operationen mit minimalem Overhead durchzuführen.

    \subsection{Umstellung auf 3D: \texttt{Vec3} und Möller-Trumbore-Algorithmus}
    \label{subsec:3d-migration}

    Die aktuelle Implementierung der Physik-Engine arbeitet vollständig im
    dreidimensionalen Raum. Im Vergleich zur ursprünglichen zweidimensionalen
    Grundidee mussten dafür die zentralen geometrischen Datenstrukturen und die
    Kollisionsberechnung angepasst werden.

    Ziel dieser Umstellung war es, Schallstrahlen nicht mehr ausschließlich in
    einer Ebene zu simulieren, sondern deren Ausbreitung in einem Raum mit Breite,
    Tiefe und Höhe zu ermöglichen.

    Positionen und Richtungen werden daher nicht mehr mit zwei Koordinaten,
    sondern mit dem Datentyp \texttt{Vec3} beschrieben. Ein Punkt im Raum besitzt
    eine \(x\)-, \(y\)- und \(z\)-Koordinate:

    \[
        P = (x,y,z)
    \]

    Diese Erweiterung betrifft insbesondere die Position des Lautsprechers, die
    Position des Mikrofons, die Startpunkte der Rays, deren Richtungsvektoren
    sowie die Reflexionspunkte an den Raumflächen.

    \subsubsection{Neue 3D-Datenstrukturen: \texttt{Ray} und \texttt{Triangle}}
    \label{subsubsec:ray-triangle}

    Die Grundlage der Umstellung war die Anpassung der zentralen Datenstrukturen.
    Ein Schallstrahl wird in der neuen Implementierung durch die Struktur
    \texttt{Ray} beschrieben. Diese enthält den aktuellen Ursprung des Strahls und
    seine Richtung im dreidimensionalen Raum:

    \begin{lstlisting}[language=Rust]
    pub struct Ray {
        pub origin: Vec3,
        pub direction: Vec3,
    }
    \end{lstlisting}

    Der Wert \texttt{origin} beschreibt den aktuellen Startpunkt des Strahls.
    Zu Beginn entspricht dieser Punkt der Lautsprecherposition. Nach jeder
    Reflexion wird der berechnete Trefferpunkt an einer Fläche zum neuen Ursprung
    des reflektierten Strahls.

    Der Wert \texttt{direction} beschreibt die Ausbreitungsrichtung des Strahls.
    Dieser Richtungsvektor wird normalisiert verwendet, damit nur die Richtung und
    nicht die Länge des Vektors eine Rolle spielt.

    Auch die Modellierung der Raumflächen musste verändert werden. In der
    zweidimensionalen Version konnten Wände als Liniensegmente dargestellt werden.
    In einem dreidimensionalen Raum muss ein Strahl jedoch mit Flächen kollidieren.
    Aus diesem Grund werden Raumflächen nun als Dreiecke modelliert:

    \begin{lstlisting}[language=Rust]
    pub struct Triangle {
        pub v0: Vec3,
        pub v1: Vec3,
        pub v2: Vec3,
        pub absorption: f32,
    }
    \end{lstlisting}

    Ein Dreieck besteht aus den drei Eckpunkten \texttt{v0}, \texttt{v1} und
    \texttt{v2}. Diese Punkte definieren eindeutig eine Fläche im
    dreidimensionalen Raum.

    Zusätzlich besitzt jedes Dreieck einen Absorptionswert. Dieser beschreibt, wie
    stark ein Schallstrahl bei einer Reflexion an der jeweiligen Fläche
    abgeschwächt wird. Damit enthält die Struktur neben der geometrischen
    Beschreibung auch eine akustische Eigenschaft.

    \subsubsection{Ablösung des 2D-Segment-Schnitts}
    \label{subsubsec:replace-2d-intersection}

    Der vorherige Schnittpunktalgorithmus war auf zweidimensionale
    Strahl-Segment-Schnitte ausgelegt. Dabei wurde geprüft, ob ein Ray eine
    Wandlinie innerhalb eines begrenzten Segments trifft.

    Dieses Verfahren ist für einen dreidimensionalen Raum nicht ausreichend, da
    Wände, Boden und Decke nicht als Linien, sondern als Flächen betrachtet werden
    müssen.

    Für die 3D-Version wurde daher auf Strahl-Dreieck-Schnittpunkte umgestellt.
    Dreiecke eignen sich besonders für diese Aufgabe, da sie immer planar sind und
    komplexe Flächen aus mehreren Dreiecken zusammengesetzt werden können.

    Dadurch kann die Simulation nicht nur einfache rechteckige Räume, sondern
    grundsätzlich auch komplexere polygonale Geometrien verarbeiten.

    \subsubsection{Aufbau des 3D-Raums aus Dreiecken}
    \label{subsubsec:room-triangles}

    Obwohl der Schnittpunktalgorithmus mit Dreiecken arbeitet, bestehen einfache
    Raumflächen wie Boden, Decke oder Wände meistens aus Rechtecken. Deshalb wird
    jede rechteckige Raumfläche in zwei Dreiecke zerlegt.

    Dies geschieht in der Hilfsfunktion
    \texttt{add\_rect\_as\_two\_triangles}:

    \begin{lstlisting}[language=Rust]
    fn add_rect_as_two_triangles(
        triangles: &mut Vec<Triangle>,
        a: Vec3,
        b: Vec3,
        c: Vec3,
        d: Vec3,
        absorption: f32,
    ) {
        triangles.push(Triangle {
            v0: a,
            v1: b,
            v2: c,
            absorption,
        });

        triangles.push(Triangle {
            v0: a,
            v1: c,
            v2: d,
            absorption,
        });
    }
    \end{lstlisting}

    Das erste Dreieck besteht aus den Punkten \texttt{a}, \texttt{b} und
    \texttt{c}. Das zweite Dreieck besteht aus den Punkten \texttt{a},
    \texttt{c} und \texttt{d}. Zusammen bilden beide Dreiecke wieder die
    vollständige rechteckige Fläche.

    Dadurch kann der Möller-Trumbore-Algorithmus einheitlich auf alle Raumflächen
    angewendet werden.

    Im Hauptprogramm wird zunächst eine Liste für alle Dreiecke des Raums
    angelegt:

    \begin{lstlisting}[language=Rust]
    let mut room: Vec<Triangle> = Vec::new();
    \end{lstlisting}

    Anschließend werden Boden, Decke und die vier Wände des Testraums
    hinzugefügt. Der Raum besitzt im Testaufbau eine Breite von \(10{,}0\), eine
    Tiefe von \(10{,}0\) und eine Höhe von \(3{,}0\).

    Das Koordinatensystem ist so aufgebaut, dass \(x\) die horizontale Richtung,
    \(y\) die Tiefe des Raums und \(z\) die Höhe beschreibt. Ein Beispiel ist der
    Boden des Raums:

    \begin{lstlisting}[language=Rust]
    add_rect_as_two_triangles(
        &mut room,
        Vec3::new(0.0, 0.0, 0.0),
        Vec3::new(width, 0.0, 0.0),
        Vec3::new(width, depth, 0.0),
        Vec3::new(0.0, depth, 0.0),
        absorption,
    );
    \end{lstlisting}

    Alle vier Punkte besitzen hier die \(z\)-Koordinate \(0{,}0\), da sich der
    Boden auf der Höhe null befindet. Für die Decke wird dagegen die
    \(z\)-Koordinate \texttt{height} verwendet.

    Bei den Wänden bleibt jeweils eine Koordinate konstant, während die beiden
    anderen Koordinaten die jeweilige Wandfläche aufspannen.

    \subsubsection{Möller-Trumbore-Algorithmus}
    \label{subsubsec:moller-trumbore}

    Zur Berechnung von Strahl-Dreieck-Schnittpunkten wurde der
    Möller-Trumbore-Algorithmus implementiert. Dieser Algorithmus wird häufig im
    Raytracing verwendet, da er direkt prüft, ob ein Ray ein Dreieck trifft, ohne
    zuvor explizit die Ebenengleichung des Dreiecks bestimmen zu müssen.

    Die Funktion \texttt{check\_intersection} bildet den Kern dieser
    Schnittpunktprüfung. Sie erhält einen Ray und ein Dreieck. Falls ein gültiger
    Treffer existiert, gibt sie einen neuen reflektierten Ray zurück. Liegt kein
    gültiger Treffer vor, wird \texttt{None} zurückgegeben.

    Zunächst werden aus den drei Eckpunkten des Dreiecks zwei Kantenvektoren
    berechnet:

    \begin{lstlisting}[language=Rust]
    let edge1 = triangle.v1 - triangle.v0;
    let edge2 = triangle.v2 - triangle.v0;
    \end{lstlisting}

    Diese beiden Kantenvektoren spannen die Dreiecksfläche im Raum auf.
    Anschließend wird ein Kreuzprodukt zwischen der Strahlrichtung und der zweiten
    Dreieckskante gebildet:

    \begin{lstlisting}[language=Rust]
    let h = ray.direction.cross(edge2);
    let a = edge1.dot(h);
    \end{lstlisting}

    Der Wert \texttt{a} zeigt an, ob der Ray parallel zur Dreiecksfläche
    verläuft. Ist dieser Wert nahe null, kann kein eindeutiger Schnittpunkt
    berechnet werden. Deshalb wird mit einem kleinen Epsilon-Wert gearbeitet:

    \begin{lstlisting}[language=Rust]
    let epsilon = 1e-6;

    if a.abs() < epsilon {
        return None;
    }
    \end{lstlisting}

    Diese Toleranz ist notwendig, da Gleitkommazahlen nicht immer exakt berechnet
    werden. Ohne den Schwellenwert könnten numerische Rundungsfehler zu falschen
    Treffern führen.

    Anschließend wird die erste baryzentrische Koordinate \(u\) berechnet:

    \begin{lstlisting}[language=Rust]
    let f = 1.0 / a;
    let s = ray.origin - triangle.v0;
    let u = f * s.dot(h);

    if u < 0.0 || u > 1.0 {
        return None;
    }
    \end{lstlisting}

    Liegt \(u\) außerhalb des Bereichs von \(0{,}0\) bis \(1{,}0\), trifft der
    Ray möglicherweise die Ebene des Dreiecks, aber nicht die begrenzte
    Dreiecksfläche selbst.

    Danach wird die zweite baryzentrische Koordinate \(v\) berechnet:

    \begin{lstlisting}[language=Rust]
    let q = s.cross(edge1);
    let v = f * ray.direction.dot(q);

    if v < 0.0 || u + v > 1.0 {
        return None;
    }
    \end{lstlisting}

    Die Bedingung \(u + v > 1{,}0\) stellt sicher, dass sich der Trefferpunkt
    nicht außerhalb des Dreiecks befindet.

    Ein gültiger Treffer liegt nur dann vor, wenn die baryzentrischen Koordinaten
    innerhalb der Dreiecksgrenzen bleiben.

    Zum Schluss wird der Parameter \(t\) berechnet. Dieser beschreibt die
    Entfernung des Schnittpunkts entlang der Strahlrichtung:

    \begin{lstlisting}[language=Rust]
    let t = f * edge2.dot(q);

    if t < 0.001 {
        return None;
    }
    \end{lstlisting}

    Negative oder sehr kleine Werte werden verworfen. Dadurch werden Treffer
    ignoriert, die hinter dem Ursprung des Rays liegen oder unmittelbar am
    Startpunkt auftreten.

    Diese Prüfung verhindert sogenannte Selbsttreffer, bei denen ein
    reflektierter Ray direkt wieder dieselbe Fläche trifft.

    \subsubsection{Reflexionsberechnung}
    \label{subsubsec:3d-reflection}

    Die zugrunde liegende Reflexionsgleichung entspricht der bereits in
    Abschnitt~\ref{subsec:reflection-vector-calculation} beschriebenen
    Vektorreflexion.

    Im dreidimensionalen Fall wird die Oberflächennormale jedoch nicht aus einem
    zweidimensionalen Wandvektor, sondern aus den beiden Kantenvektoren des
    getroffenen Dreiecks berechnet.

    Sobald ein gültiger Schnittpunkt gefunden wurde, wird zunächst der tatsächliche
    Trefferpunkt im dreidimensionalen Raum bestimmt:

    \begin{lstlisting}[language=Rust]
    let hit_point = ray.origin + ray.direction * t;
    \end{lstlisting}

    Dieser Punkt wird anschließend zum neuen Ursprung des reflektierten Rays.

    Danach wird die Oberflächennormale des Dreiecks berechnet. Die Normale steht
    senkrecht auf der Dreiecksfläche und ergibt sich aus dem Kreuzprodukt der
    beiden Kantenvektoren:

    \begin{lstlisting}[language=Rust]
    let normal = edge1.cross(edge2).normalize();
    \end{lstlisting}

    Die neue Richtung des Schallstrahls wird berechnet, indem die ursprüngliche
    Richtung an dieser Normalen gespiegelt wird:

    \begin{lstlisting}[language=Rust]
    let bounced_direction =
        ray.direction.reflect(normal).normalize();
    \end{lstlisting}

    Die Methode \texttt{reflect} des \texttt{glam}-Crates übernimmt die
    mathematische Spiegelung des Richtungsvektors. Anschließend wird die neue
    Richtung normalisiert, damit sie weiterhin als Einheitsvektor verwendet werden
    kann.

    Die Funktion gibt schließlich einen neuen Ray zurück:

    \begin{lstlisting}[language=Rust]
    Some(Ray {
        origin: hit_point,
        direction: bounced_direction,
    })
    \end{lstlisting}

    Der reflektierte Ray startet somit am berechneten Trefferpunkt und bewegt sich
    in der neuen Reflexionsrichtung weiter. Dadurch können mehrere Reflexionen
    eines Schallstrahls nacheinander simuliert werden.

    \subsubsection{Mikrofonerkennung im 3D-Raum}
    \label{subsubsec:3d-microphone}

    Neben der Reflexion an Dreiecksflächen musste auch die Erkennung eines
    Mikrofontreffers an den dreidimensionalen Raum angepasst werden.

    Das Mikrofon wird in der Simulation als Kugel mit einem konfigurierbaren
    Radius betrachtet. Ein Ray muss daher nicht exakt den Mittelpunkt des
    Mikrofons treffen, sondern nur nahe genug an diesem vorbeifliegen.

    Die Funktion \texttt{check\_mic\_intersection} prüft hierfür den Abstand
    zwischen einem Ray-Segment und dem Mikrofonmittelpunkt.

    Wichtig ist, dass nicht der unendliche Ray betrachtet wird, sondern nur das
    konkrete Segment zwischen dem Startpunkt des aktuellen Rays und dem nächsten
    Trefferpunkt an einer Fläche:

    \begin{lstlisting}[language=Rust]
    pub fn check_mic_intersection(
        start: Vec3,
        end: Vec3,
        mic_center: Vec3,
        mic_radius: f32,
    ) -> bool {
        let segment = end - start;
        let segment_length_sq =
            segment.length_squared();

        if segment_length_sq == 0.0 {
            return start.distance(mic_center)
                <= mic_radius;
        }

        let to_mic = mic_center - start;

        let t =
            to_mic.dot(segment)
            / segment_length_sq;

        let clamped_t = t.clamp(0.0, 1.0);

        let closest_point =
            start + segment * clamped_t;

        closest_point.distance(mic_center)
            <= mic_radius
    }
    \end{lstlisting}

    Zunächst wird der Segmentvektor zwischen Start- und Endpunkt berechnet.
    Danach wird der Mikrofonmittelpunkt auf dieses Segment projiziert.

    Der berechnete Projektionswert wird mit
    \texttt{clamp(0.0, 1.0)} begrenzt. Dadurch wird sichergestellt, dass nur das
    tatsächliche Segment und nicht die unendliche Verlängerung des Rays betrachtet
    wird.

    Anschließend wird der Punkt auf dem Segment berechnet, der dem Mikrofon am
    nächsten liegt. Ist der Abstand zwischen diesem Punkt und dem
    Mikrofonmittelpunkt kleiner oder gleich dem Mikrofonradius, wird ein Treffer
    registriert.

    Diese Methode ist robuster als eine reine Punktprüfung, da es bei zufällig
    ausgesendeten Rays sehr unwahrscheinlich wäre, den exakten Mittelpunkt des
    Mikrofons zu treffen.

    \subsubsection{Integration in die Simulationsschleife}
    \label{subsubsec:simulation-loop}

    Die neue Schnittpunktprüfung wird in der Simulationslogik verwendet, um jeden
    einzelnen Schallstrahl über mehrere Reflexionen hinweg zu verfolgen.

    Zunächst wird ein zufälliger Startstrahl an der Lautsprecherposition erzeugt:

    \begin{lstlisting}[language=Rust]
    fn generate_initial_ray(
        config: &SimulationConfig
    ) -> Ray {
        let mut rng = rand::rng();

        let theta: f32 =
            rng.random_range(0.0..(2.0 * PI));

        let z: f32 =
            rng.random_range(-1.0..1.0);

        let r =
            (1.0_f32 - z * z).sqrt();

        let direction = Vec3::new(
            r * theta.cos(),
            r * theta.sin(),
            z,
        ).normalize();

        Ray {
            origin: config.speaker_position,
            direction,
        }
    }
    \end{lstlisting}

    Anders als in einer zweidimensionalen Simulation reicht hier kein einzelner
    Winkel in einer Ebene aus. Stattdessen wird eine zufällige Richtung auf der
    Einheitskugel erzeugt.

    Der Winkel \(\theta\) beschreibt die Richtung um die \(z\)-Achse. Der Wert
    \(z\) bestimmt den Anteil der Richtung nach oben oder unten. Aus diesen
    Werten entsteht ein normalisierter \texttt{Vec3}-Richtungsvektor.

    Für die eigentliche Kollision wird in der Funktion \texttt{cast\_ray} das
    nächstgelegene getroffene Dreieck gesucht:

    \begin{lstlisting}[language=Rust]
    for triangle in triangles {
        if let Some(bounced_ray) =
            geometry::check_intersection(
                ray,
                triangle
            )
        {
            let distance =
                ray.origin.distance(
                    bounced_ray.origin
                );

            if distance < shortest_distance {
                shortest_distance = distance;

                closest_bounced_ray = Some((
                    bounced_ray,
                    triangle.absorption
                ));
            }
        }
    }
    \end{lstlisting}

    Dabei wird jedes Dreieck des Raums geprüft. Wenn mehrere gültige Treffer
    existieren, wird der Treffer mit der kürzesten Distanz ausgewählt.

    Dies ist notwendig, da ein Ray zuerst auf die nächstgelegene Fläche treffen
    muss und nicht auf eine weiter entfernte Fläche dahinter.

    In der Funktion \texttt{simulate\_single\_ray} wird dieser Ablauf über
    mehrere Reflexionen wiederholt. Nach jedem Treffer wird geprüft, ob der Ray auf
    seinem Flugsegment am Mikrofon vorbeigeflogen ist:

    \begin{lstlisting}[language=Rust]
    if geometry::check_mic_intersection(
        start_point,
        end_point,
        config.mic_position,
        config.mic_radius,
    ) {
        let distance_to_mic =
            start_point.distance(
                config.mic_position
            );

        let total_distance_at_hit =
            total_distance + distance_to_mic;

        if total_distance_at_hit > 0.001 {
            ray_hits.push((
                total_distance_at_hit / 343.0,
                current_pressure,
            ));
        }
    }
    \end{lstlisting}

    Wird ein Mikrofontreffer erkannt, werden die Ankunftszeit und der aktuelle
    Druckwert gespeichert.

    Die Ankunftszeit wird in der aktuellen Implementierung näherungsweise über die
    zurückgelegte Strecke berechnet und anschließend durch die
    Schallgeschwindigkeit geteilt. Dafür wird der Näherungswert
    \(343{,}0\,\mathrm{m/s}\) verwendet.

    Entsprechend dem in Abschnitt~\ref{subsec:akustik} beschriebenen
    Absorptionsmodell wird der aktuelle Druckwert nach jeder Reflexion reduziert:

    \begin{lstlisting}[language=Rust]
    current_pressure *= 1.0 - absorption;

    if current_pressure < config.min_pressure {
        break;
    }
    \end{lstlisting}

    Der Absorptionswert der getroffenen Fläche bestimmt, wie stark der Ray nach
    der Reflexion abgeschwächt wird.

    Sinkt der Druck unter den konfigurierten Mindestwert, wird die Simulation
    dieses Rays beendet. Dadurch werden sehr leise und für das Ergebnis kaum
    relevante Reflexionen nicht unnötig weiterberechnet.

    \subsubsection{Konfiguration und Visualisierung der 3D-Simulation}
    \label{subsubsec:3d-config-visualization}

    Die Parameter der Simulation werden aus der Datei
    \texttt{input\_config.json} geladen.

    Die zugehörige Struktur \texttt{SimulationConfig} enthält unter anderem die
    maximale Anzahl an Reflexionen, den minimalen Druckwert, den Mikrofonradius,
    die Anzahl der auszusendenden Rays sowie die Positionen von Lautsprecher und
    Mikrofon:

    \begin{lstlisting}[language=Rust]
    #[derive(Deserialize)]
    pub struct SimulationConfig {
        pub max_bounces: u32,
        pub min_pressure: f32,
        pub mic_radius: f32,
        pub mic_position: Vec3,
        pub rays_to_cast: u32,
        pub speaker_position: Vec3,
    }
    \end{lstlisting}

    Auch an dieser Stelle zeigt sich die Umstellung auf den dreidimensionalen
    Raum: Sowohl \texttt{mic\_position} als auch
    \texttt{speaker\_position} sind nun vom Typ \texttt{Vec3}.

    In der Konfigurationsdatei werden diese Positionen deshalb als Arrays mit drei
    Zahlen angegeben:

    \begin{lstlisting}
    {
      "max_bounces": 32,
      "min_pressure": 0.001,
      "mic_radius": 1.5,
      "mic_position": [5.0, 5.0, 1.5],
      "rays_to_cast": 1000000,
      "speaker_position": [2.0, 2.0, 1.5]
    }
    \end{lstlisting}

    Die dritte Koordinate beschreibt die Höhe im Raum.

    Die Lautsprecherposition \texttt{[2.0, 2.0, 1.5]} bedeutet beispielsweise,
    dass der Lautsprecher bei \(x = 2{,}0\), \(y = 2{,}0\) und einer Höhe von
    \(z = 1{,}5\) Metern platziert ist.

    Gleiches gilt für die Mikrofonposition. Dadurch kann die Simulation die
    Schallausbreitung nicht nur auf einer zweidimensionalen Ebene, sondern
    innerhalb eines vollständigen 3D-Raums berechnen.

    Neben der Konfiguration wurden auch die Visualisierungsdaten an die
    dreidimensionale Darstellung angepasst.

    Die Struktur \texttt{VisualizerOutput} speichert Lautsprecher, Mikrofon und
    die Punkte der Ray-Pfade ebenfalls als \texttt{Vec3}:

    \begin{lstlisting}[language=Rust]
    pub struct VisualizerOutput {
        pub speaker: Vec3,
        pub mic: Vec3,
        pub mic_radius: f32,
        pub rays: Vec<Vec<Vec3>>,
    }
    \end{lstlisting}

    Die Variable \texttt{rays} enthält mehrere Ray-Pfade. Jeder einzelne Pfad
    besteht aus mehreren Punkten im Raum, wobei jeder Punkt durch drei Koordinaten
    beschrieben wird.

    Dadurch können ausgewählte Strahlen mit ihren Reflexionspunkten als
    dreidimensionale Pfade ausgegeben werden.

    Die Visualisierungsdaten dienen vor allem zur Kontrolle und zum Debugging. Sie
    ermöglichen zu prüfen, ob die Rays korrekt von der Lautsprecherposition
    ausgesendet werden, ob sie an den Dreiecksflächen reflektieren und ob die
    berechneten Pfade im Raum plausibel verlaufen.

    Damit unterstützt die Visualisierung die Überprüfung der neuen
    3D-Kollisions- und Reflexionslogik.

    \subsubsection{Ergebnis der Umstellung}
    \label{subsubsec:3d-result}

    Durch die Migration von \texttt{Vec2} auf \texttt{Vec3} wurde die bisherige
    Beschränkung auf eine zweidimensionale Simulation aufgehoben.

    Die Engine kann nun Schallstrahlen in einem dreidimensionalen Raum erzeugen,
    deren Kollisionen mit Dreiecksflächen berechnen und die reflektierten Rays
    weiterverfolgen.

    Die Einführung des Möller-Trumbore-Algorithmus ermöglicht eine robuste und
    direkte Prüfung von Strahl-Dreieck-Schnittpunkten.

    In Kombination mit der Reflexionsberechnung über Dreiecksnormalen, der
    Mikrofonerkennung als Segment-Kugel-Prüfung und dem Export der
    dreidimensionalen Ray-Pfade bildet diese Umstellung die Grundlage für eine
    realistischere akustische Simulation.

    Gleichzeitig bleibt die Architektur modular, da Geometrie, Simulation,
    Konfiguration und Export klar voneinander getrennt sind.


    \section{Hochleistungs-Nebenläufigkeit und Datenparallelität}
    \label{sec:concurrency}

    %A primary advantage of Rust in high-performance computing is its strict compile-time guarantees against data races, colloquially known as
    %``fearless concurrency.'' Distributing a heavily iterative physics loop across multi-core processor architectures in traditional languages often
    %requires complex thread-locking logic (e.g., Mutexes) to prevent memory corruption.
    %This engine avoids locking overhead entirely by employing a functional, data-parallel paradigm utilizing the \texttt{rayon} crate.

    Ein wesentlicher Vorteil von Rust im High-Performance-Computing sind seine strengen Garantien zur Kompilierzeit gegen Data-Races.
    Die Verteilung einer rechenintensiven, iterativen Physikschleife auf Multi-Core-Prozessorarchitekturen erfordert in traditionellen
    Programmiersprachen oft komplexe Thread-Sperrlogiken (z.B. Mutexe), um Speicherbeschädigungen zu verhindern.
    In dieser Engine wird dieser Overhead vermieden durch ein funktionales, datenparalleles Paradigma unter Verwendung des \texttt{rayon}-Crates.

    \subsection{Architektonischer Wandel: Von Shared Mutation zu thread-lokalem-Folding}\label{subsec:architectural-shift:-from-shared-mutation-to-thread-local-folding}

    %In a standard single-threaded execution, sequential iterations append calculated intersection data (e.g., delay and pressure values) into a shared dynamic array.
    %Forcing this loop across numerous threads would result in concurrent write attempts to the same memory addresses, leading to immediate data races.
    %Rather than wrapping these shared arrays in computationally expensive, thread-safe locks, the architecture leverages \texttt{rayon}'s \texttt{into\_par\_iter()} trait
    %combined with a \texttt{fold} and \texttt{reduce} pipeline.

    Bei einer standardmäßigen Single-Thread-Ausführung hängen sequenzielle Iterationen die berechneten Schnittpunktdaten (z.B. delay und pressure werte) an ein gemeinsames dynamisches Array an.
    Diese Schleife auf zahlreiche Threads aufzuteilen, würde zur gleichzeitigen Schreibversuchen auf dieselben Speicheradressen führen, was unmittelbare Data Races zur Folge hätte.
    Anstatt diese gemeinsam genutzt Arrays in rechenintensive, threadsichere Sperren (Locks) zu hüllen, nutzt die Architektur das Trait \texttt{into\_par\_iter()} von \texttt{rayon} in
    Kombination mit einer \texttt{fold}- und \texttt{reduce}- Pipeline.

    Dieser Ansatz teilt die Arbeitslast sicher auf alle verfügbaren CPU-Kerne auf.
    Anstatt zu versuchen, in ein globales Array zu schreiben oder Speicher für jeden Strahl
    dynamisch zu allokieren, verwendet die Engine die \texttt{fold}-Operation, um einen privaten, vorab allokieren, thread-lokalen Puffer für jeden Kern bereitzustellen.

    Jeder Thread simuliert unabhängig seine zugewiesene Teilmenge an Strahlen und hängt die Schnittpunktergebnisse sequenziell in seinen eigenen lokalen Speicherbereich an, ohne Sperren (lock-free).
    Sobald ein Thread seine Arbeitslast abgeschlossen hat, wird die \texttt{reduce}-Phase aktiviert, welche diese isolierten thread-lokalen Puffer sicher und effizient zu den endgültigen,
    vereinheitlichten Arrays zusammenführt.
    Dieser funktionale Ansatz garantiert die Datensicherheit und maximiert gleichzeitig die Hardwareauslastung.

    %This approach safely divides the workload—the generation and simulation of acoustic rays—among all available CPU cores.
    %Instead of attempting to write to a global array or dynamically allocating memory for every individual ray, the engine utilizes the \texttt{fold} operation to provision a private,
    %pre-allocated, thread-local buffer for each core.

    %Each thread independently simulates its assigned subset of rays and sequentially appends the intersection results into its own local memory space, completely lock-free.
    %Once a thread completes its workload, the \texttt{reduce} phase activates, safely and efficiently merging these isolated thread-local buffers into the final, unified arrays.
    %This functional approach guarantees data safety while maximizing hardware utilization.

    \begin{lstlisting}[language=Rust]
use rayon::prelude::*;

fn run_simulation_parallel(
config: &SimulationConfig,
walls: &Vec<Wall>
) -> (Vec<f32>, Vec<f32>) {
    let(delays_par, pressures_par, _) = (0..config.rays_to_cast)
        .into_par_iter()
        // 1. Falten, Split
        .fold( || {
            let local_delays = Vec::with_capacity(CAPACITY);
            let local_pressures = Vec::with_capacity(CAPACITY);
            let temp_ray_buffer = Vec::with_capacity(CAP);

            (local_delays,local_pressures,temp_ray_buffer)
        },
            |mut local_buf, _| {
                let fresh_ray = generate_initial_ray(config);
                simulate_single_ray(fresh_ray,walls,config,&mut local_buf.2);
                for(delay, pressure) in &local_buf.2 {
                    local_buf.0.push(*delay);
                    local_buf.1.push(*pressure);
                }
                local_buf
            }
        // 2. Reduzieren, Global Merge
        ).reduce(
        || (Vec::new(),Vec::new(),Vec::new()),
        |mut a, mut b| {
            a.0.append(&mut b.0);
            a.1.append(&mut b.1);
            a
        }
    );
    (delays_par, pressures_par)
}
    \end{lstlisting}


    \newpage


    \section{Systemarchitektur und Pipeline-Orchestrierung}
    \label{sec:systemarchitektur}

    Die Rust-Physik-Engine bildet den rechenintensiven Kern der Simulation.
    Die weitere Verarbeitung der erzeugten Impulsantwort erfolgt in der
    Python-Pipeline.

    Um beide Teile miteinander zu verbinden, wurde eine Schnittstelle
    implementiert, die die kompilierte Rust-Binary aus Python heraus startet und
    anschließend die erzeugte Datei \texttt{ir\_output.json} einliest.

    Die Schnittstelle befindet sich in der Datei
    \texttt{python\_service/rust\_runner.py}. Ihre zentrale Funktion ist
    \texttt{run\_rust\_and\_load\_json}.

    Sie übernimmt die Orchestrierung zwischen Python und Rust: Python startet den
    Rust-Prozess, wartet auf dessen Abschluss, prüft die erzeugte Ausgabedatei und
    stellt die geladenen Daten anschließend als Python-Dictionary bereit.

    \[
        \text{Python-Pipeline}
        \rightarrow
        \texttt{rust\_runner.py}
        \rightarrow
        \text{Rust-Binary}
        \rightarrow
        \texttt{ir\_output.json}
        \rightarrow
        \text{Python-Dictionary}
    \]

    Diese Architektur trennt die Zuständigkeiten klar voneinander.

    Rust bleibt für die performante Berechnung der akustischen Simulation
    verantwortlich, während Python die Ausführung koordiniert und die
    Ergebnisdaten für die spätere Audioverarbeitung vorbereitet.

    \subsection{Aufgabe der Python-Rust-Schnittstelle}
    \label{subsec:python-rust-interface}

    Die Funktion \texttt{run\_rust\_and\_load\_json} erhält zwei Parameter:
    den Pfad zur kompilierten Rust-Binary und den erwarteten Namen der
    JSON-Ausgabedatei.

    Beide Werte werden zunächst in \texttt{Path}-Objekte umgewandelt, damit Pfade
    plattformunabhängig und zuverlässig verarbeitet werden können.

    \begin{lstlisting}[language=Python]
    def run_rust_and_load_json(
        rust_binary_path: str,
        json_path: str
    ) -> dict:
        rust_binary = Path(
            rust_binary_path
        ).resolve()

        json_file = Path(json_path)
    \end{lstlisting}

    Der Rückgabetyp der Funktion ist ein \texttt{dict}. Dadurch stehen die von
    Rust erzeugten Daten nach erfolgreicher Ausführung direkt als
    Python-Dictionary zur Verfügung.

    Spätere Pipeline-Schritte müssen die JSON-Datei daher nicht erneut selbst
    öffnen und parsen.

    Die Schnittstelle erfüllt drei zentrale Aufgaben:

    \begin{itemize}
        \item Start der Rust-Binary als externer Prozess,
        \item Prüfung und Validierung der Ausführung,
        \item Einlesen der erzeugten Datei
        \texttt{ir\_output.json}.
    \end{itemize}

    \subsection{Pfadprüfung und Arbeitsverzeichnis}
    \label{subsec:path-validation}

    Vor dem Start der Rust-Binary wird geprüft, ob der übergebene Pfad existiert
    und ob es sich tatsächlich um eine Datei handelt.

    Dadurch werden ungültige Pfadkonfigurationen bereits vor dem Start des
    Subprozesses abgefangen.

    \begin{lstlisting}[language=Python]
    if not rust_binary.exists():
        raise FileNotFoundError(
            "Rust-Binary nicht gefunden: "
            f"{rust_binary}"
        )

    if not rust_binary.is_file():
        raise FileNotFoundError(
            "Pfad ist keine Datei: "
            f"{rust_binary}"
        )
    \end{lstlisting}

    Diese Prüfung verbessert die Robustheit der Schnittstelle.

    Ein falscher Pfad führt dadurch nicht zu einem schwer nachvollziehbaren Fehler
    innerhalb von \texttt{subprocess.run}, sondern zu einer klaren und
    verständlichen Fehlermeldung.

    Zusätzlich wird geprüft, ob als erwartete Ausgabedatei tatsächlich
    \texttt{ir\_output.json} übergeben wurde:

    \begin{lstlisting}[language=Python]
    if json_file.name != "ir_output.json":
        raise ValueError(
            "Rust exportiert aktuell fest nach "
            "'ir_output.json'. "
            "Uebergeben wurde aber: "
            f"{json_file.name}"
        )
    \end{lstlisting}

    Diese Prüfung ist notwendig, weil die Rust-Seite aktuell fest in eine Datei
    mit dem Namen \texttt{ir\_output.json} exportiert.

    Würde Python eine andere Datei erwarten, würde die Schnittstelle nach der
    Ausführung an der falschen Stelle suchen. Die Prüfung verhindert damit eine
    inkonsistente Annahme zwischen Rust und Python.

    Ein weiterer wichtiger Punkt ist das Arbeitsverzeichnis.

    Die Rust-Simulation benötigt die Datei \texttt{input\_config.json}. Da sich
    die kompilierte Binary typischerweise in einem Unterordner wie
    \texttt{target/release} befindet, wird der passende Projektordner automatisch
    gesucht:

    \begin{lstlisting}[language=Python]
    rust_service_dir = rust_binary.parent

    for parent in rust_binary.parents:
        config_file = (
            parent / "input_config.json"
        )

        if config_file.exists():
            rust_service_dir = parent
            break
    \end{lstlisting}

    Sobald ein Ordner mit \texttt{input\_config.json} gefunden wird, wird dieser
    als Arbeitsverzeichnis für den Rust-Prozess verwendet.

    Dadurch kann die Binary aus dem Build-Verzeichnis gestartet werden, während
    die Konfigurationsdatei im eigentlichen
    \texttt{rust\_service}-Ordner verbleibt.

    \subsection{Start der Rust-Binary}
    \label{subsec:start-rust-binary}

    Vor dem Start der Simulation wird eine eventuell vorhandene alte
    \texttt{ir\_output.json} gelöscht:

    \begin{lstlisting}[language=Python]
    output_file = (
        rust_service_dir / "ir_output.json"
    )

    if output_file.exists():
        output_file.unlink()
    \end{lstlisting}

    Dieser Schritt verhindert, dass nach einem fehlgeschlagenen Rust-Lauf
    versehentlich eine alte Ergebnisdatei eingelesen wird.

    Ohne diese Löschung könnte eine veraltete Datei den Eindruck erwecken, dass die
    aktuelle Simulation erfolgreich neue Daten erzeugt hat.

    Die eigentliche Ausführung der Rust-Binary erfolgt anschließend über
    \texttt{subprocess.run}:

    \begin{lstlisting}[language=Python]
    result = subprocess.run(
        [str(rust_binary)],
        cwd=rust_service_dir,
        capture_output=True,
        text=True
    )
    \end{lstlisting}

    Der Parameter \texttt{cwd} legt das Arbeitsverzeichnis des Prozesses fest.

    In diesem Fall wird der zuvor ermittelte
    \texttt{rust\_service}-Ordner verwendet, damit Rust die Datei
    \texttt{input\_config.json} findet und die Datei
    \texttt{ir\_output.json} an der erwarteten Stelle erzeugt.

    Mit \texttt{capture\_output=True} werden die Ausgaben des Rust-Prozesses
    gesammelt. Dadurch können \texttt{stdout} und \texttt{stderr} im Fehlerfall
    in die Fehlermeldung übernommen werden.

    Der Parameter \texttt{text=True} sorgt dafür, dass diese Ausgaben als Text
    und nicht als Byte-Daten verarbeitet werden.

    \subsection{Fehlerbehandlung und JSON-Ingestion}
    \label{subsec:error-handling-json}

    Nach dem Prozessstart wird zunächst der Rückgabecode der Rust-Binary geprüft.

    Ein Rückgabecode ungleich null zeigt an, dass der Rust-Prozess fehlgeschlagen
    ist:

    \begin{lstlisting}[language=Python]
    if result.returncode != 0:
        raise RuntimeError(
            "Rust-Binary ist fehlgeschlagen.\n"
            f"Returncode: "
            f"{result.returncode}\n\n"
            f"STDOUT:\n"
            f"{result.stdout}\n\n"
            f"STDERR:\n"
            f"{result.stderr}"
        )
    \end{lstlisting}

    Die Fehlermeldung enthält neben dem Returncode auch die Standardausgabe und die
    Fehlerausgabe des Rust-Prozesses.

    Das ist besonders wichtig, da die Ursache eines Fehlers häufig nicht in
    Python, sondern in der extern gestarteten Rust-Binary liegt.

    Anschließend wird geprüft, ob die erwartete Datei
    \texttt{ir\_output.json} tatsächlich erzeugt wurde:

    \begin{lstlisting}[language=Python]
    if not output_file.exists():
        raise FileNotFoundError(
            "Rust wurde ausgefuehrt, aber "
            "ir_output.json wurde nicht "
            "gefunden: "
            f"{output_file}"
        )
    \end{lstlisting}

    Diese Prüfung deckt den Fall ab, dass der Rust-Prozess zwar ohne Fehlercode
    beendet wurde, aber trotzdem keine Ausgabedatei erzeugt hat.

    Auch dieser Fall wird dadurch mit einer eindeutigen Fehlermeldung behandelt.

    Nach erfolgreicher Ausführung wird die JSON-Datei geöffnet und mit
    \texttt{json.load} in ein Python-Dictionary umgewandelt:

    \begin{lstlisting}[language=Python]
    try:
        with output_file.open(
            "r",
            encoding="utf-8"
        ) as file:
            ir_data = json.load(file)

    except json.JSONDecodeError as error:
        raise ValueError(
            "ir_output.json ist keine "
            "gueltige JSON-Datei: "
            f"{output_file}"
        ) from error

    return ir_data
    \end{lstlisting}

    Falls die Datei beschädigt ist oder kein gültiges JSON enthält, wird der
    \texttt{JSONDecodeError} abgefangen und in eine verständlichere Fehlermeldung
    umgewandelt.

    Dadurch wird verhindert, dass fehlerhafte Ergebnisdaten unbemerkt an die
    weiteren Verarbeitungsschritte weitergegeben werden.

    \subsection{Manueller Ende-zu-Ende-Test}
    \label{subsec:end-to-end-test}

    Zur Überprüfung der Schnittstelle wurde die Datei
    \texttt{test\_rust\_runner\_manual.py} erstellt.

    Dieser Test startet die Rust-Binary über die Runner-Funktion und prüft
    anschließend, ob die erzeugte JSON-Struktur korrekt geladen wurde.

    \begin{lstlisting}[language=Python]
    from rust_runner import (
        run_rust_and_load_json
    )

    def main():
        print("Testing Rust runner...")

        rust_binary = (
            "../rust_service/target/release/"
            "rust_service.exe"
        )

        json_output = "ir_output.json"

        ir_data = run_rust_and_load_json(
            rust_binary,
            json_output
        )

        print(type(ir_data))
        print(ir_data.keys())
        print(ir_data["metadata"])
        print(ir_data["hits"].keys())
    \end{lstlisting}

    Der Test kontrolliert zunächst, ob ein Python-Dictionary zurückgegeben wird.

    Danach werden die obersten Schlüssel der JSON-Struktur ausgegeben. Erwartet
    werden insbesondere die Bereiche \texttt{metadata} und \texttt{hits}.

    Innerhalb von \texttt{hits} werden die Arrays
    \texttt{delays\_seconds} und \texttt{pressures} erwartet.

    Zusätzlich werden die Längen der beiden Ergebnislisten ausgegeben:

    \begin{lstlisting}[language=Python]
    delays = (
        ir_data["hits"]["delays_seconds"]
    )

    pressures = (
        ir_data["hits"]["pressures"]
    )

    print(len(delays))
    print(len(pressures))
    \end{lstlisting}

    Dadurch kann kontrolliert werden, ob beide Listen vorhanden und plausibel
    gefüllt sind.

    Der Test dient damit als einfache Ende-zu-Ende-Prüfung der Verbindung zwischen
    Python und Rust.

    \subsection{Bewertung der Architektur}
    \label{subsec:architecture-evaluation}

    Die Kommunikation zwischen Python und Rust erfolgt bewusst über die Datei
    \texttt{ir\_output.json}.

    Diese Lösung ist nachvollziehbar, modular und benötigt keine direkten
    Rust-Python-Bindings wie beispielsweise \texttt{PyO3}. Dadurch bleiben beide
    Projektteile unabhängig voneinander ausführbar und können getrennt getestet
    werden.

    Ein wesentlicher Vorteil ist die klare Trennung der Verantwortlichkeiten:

    \begin{itemize}
        \item Rust berechnet die physikalische Simulation.
        \item Python startet und überwacht den Rust-Prozess.
        \item Python prüft die erzeugten Ausgabedaten.
        \item Python übernimmt die weitere Verarbeitung der Ergebnisse.
    \end{itemize}

    Die implementierte Fehlerbehandlung sorgt dafür, dass fehlende Dateien,
    falsche Pfade, fehlgeschlagene Rust-Prozesse und ungültige JSON-Ausgaben früh
    erkannt werden.

    Eine Einschränkung dieser Architektur besteht darin, dass die Kommunikation
    über eine Datei zusätzlichen Datei-I/O verursacht.

    Bei sehr großen Simulationen kann das Schreiben und anschließende Einlesen der
    JSON-Datei einen relevanten Anteil der Gesamtlaufzeit ausmachen. Die gemessenen
    Auswirkungen dieses Datei-I/O werden in
    Abschnitt~\ref{sec:benchmarks} genauer untersucht.

    Eine mögliche spätere Optimierung wäre die Datenübertragung über Streams oder
    eine direkte Rust-Python-Bindung.

    Für die aktuelle Projektstruktur wurde die JSON-basierte Schnittstelle jedoch
    bewusst gewählt, da sie transparent, robust und einfach testbar ist.

    Damit bildet \texttt{rust\_runner.py} die technische Verbindung zwischen dem
    Rust-Kern und der späteren Python-basierten Audioverarbeitung.


    \section{Leistungsoptimierungen und Profiling}
    \label{sec:optimizations}
    Benchmarking the compiled binary (utilizing Rust's \texttt{--release} flag for maximum LLVM optimization)
    revealed significant bottlenecks and subsequent optimization opportunities within the parallelized architecture.

    Das Benchmarking der kompilierten Binärdatei (unter Verwendung des \texttt{--release} flag und dem \texttt{let var = Instant::now();} timer) deckte signifikante
    Engpässe und darauffolgende Optimierungsmöglichkeiten innerhalb der parallelisierten Architektur auf.

    \subsection{Reduzierung des Speicherallokations-Overheads: \texttt{flat\_map} vs. \texttt{fold}}
    \label{sec:memory_overhead}


    Während die anfängliche parallele Implementierung den \texttt{flat\_map}-iterator von \texttt{rayon} nutzt, offenbarte die Skalierung der Simulation auf $100.000.000$ Strahlen einen
    signifikanten Leistungsengpass bezüglich der Speicherallokation.
    Die \texttt{flat\_map}-Operation allozierte dynamisch einen neuen \texttt{Vec} für die Rückgabedaten jedes einzelnen Strahls.
    Bei der Verteilung auf alle CPU-Kerne führte dies zu Millionen von flüchtigen Heap-Allokationen pro Sekunde.
    Das System erlebte Memory Thrashing; die CPU-Auslastung fiel auf 0\% der Hauptspeicher war ausgelastet, und das Betriebssystem entzog dem Ausführungs-Thread die Ressourcen,
    was zum Stillstand der Simulation führte.

    Um diesen Allokations-Overhead zu beheben, wurde die parallele Pipeline refaktorisiert, um das \texttt{fold}-Muster zu verwenden.
    Anstatt Speicher pro Strahl zu allokieren, stellt die \texttt{fold}-Operation exakt einen thread-lokalen Puffer (initialisiert über \texttt{Vec::with\_capacity();}) pro CPU-Kern bereit.
    Die Kerne akkumulieren ihre lokalen Ergebnisse und führen diese erst nach Abschluss des Threads zusammen.
    Diese architektonische Änderung reduzierte die Heap-Allokationen um mehrere Größenordnungen und ermöglichte es dem System, $100.000.000$ Strahlen effizient zu verarbeiten,
    ohne die Speicherressourcen zu erschöpfen.
    %While the initial parallel implementation utilized \texttt{rayon}'s \texttt{flat\_map} iterator, scaling the simulation to $100,000,000$ rays exposed a critical flaw in memory allocation.
    %The \texttt{flat\_map} operation dynamically allocated a new \texttt{Vec} for every single ray's return data.
    %When distributed across all CPU cores, this resulted in millions of transient heap allocations per second.
    %The system experienced severe memory thrashing; CPU utilization dropped to 0\%, system memory was exhausted, and the OS starved the execution thread, causing the simulation to stall.
    %
    %To resolve this allocator crisis, the parallel pipeline was refactored to use the \texttt{fold} pattern.
    %Instead of allocating memory per ray, the \texttt{fold} operation provisions exactly one thread-local buffer (initialized via \texttt{Vec::with\_capacity()}) per CPU core.
    %The cores accumulate their local results and merge them only upon thread completion.
    %This architectural shift reduced heap allocations by orders of magnitude, allowing the system to process $100,000,000$ rays without memory exhaustion.

    \subsection{Benchmarking Ergebnisse und I/O Overhead}
    \label{sec:benchmarks}

    Empirische Tests verdeutlichen die Wirksamkeit der optimierten Parallelisierung sowie die verborgenen Kosten von System-I/O und Thread-Initialisierung.
    %Empirical testing highlighted the effectiveness of the optimized parallelization, as well as the hidden costs of system I/O and thread initialization.

    \begin{itemize}
        \item \textbf{Paralleler Speedup:} Bei einer moderaten Arbeitslast wurde die Single-Thread-Simulation in $11,79$ Sekunden abgeschlossen.
        Die durch \texttt{rayon} optimierte parallele Simulation bewältigt dieselbe Arbeitslast in $2,56$ Sekunden, was einen $4{,}60\times$ Speedup auf der Test-Hardware ergab.
        \item \textbf{Thread-Pool-Kaltstart:} Bei extrem kleinen Arbeitslasten wies \texttt{rayon} einen Initialisierung-Overhead auf.
        In einem Benchmark mit wenigen Strahlen wurde die Single-Thread-Ausführung in $0,14$ Sekunden abgeschlossen, während die parallele Ausführung $0,34$ Sekunden in Anspruch nahm.
        Dies deutet darauf hin, dass die Kosten für die Initialisierung des Thread-Pools die Recheneinsparungen bei trivialen Arbeitslasten überwiegen
        \item \textbf{The I/O Engpass:} Das Profiling der End-to-End-Ausführungszeit deckte einen massiven Engpass bei der Datenserialisierung auf.
        In einem Test mit $10.000.000$ Strahlen wurde die interne Simulation schnell ausgeführt($4,72s$ Single-Thread vs. $0,64s$ Parallel).
        Die anschließende Serialisierung und das Schreiben des Datensatzes dominierten jedoch die End-to-End-Laufzeit($37,40s$ vs. $33,24s$),
        was die Notwendigkeit eines RAM-gestützten \texttt{BufWriter} unterstreicht.
    \end{itemize}
    %\begin{itemize}
    %    \item \textbf{Parallel Speedup:} For a moderate workload, the single-threaded simulation completed in 11.79 seconds. The \texttt{rayon}-optimized parallel simulation completed the same workload in 2.56 seconds, yielding a $4.60\times$ speedup on the test hardware.
    %    \item \textbf{Thread Pool Cold Start:} For extremely small workloads, \texttt{rayon} exhibited an initialization overhead.
    %    In a low-ray benchmark, the single-threaded execution completed in 0.14 seconds, whereas the parallel execution took 0.34 seconds.
    %    This indicates that the cost of initializing the thread pool outweighs the computational savings for trivial workloads.
    %    \item \textbf{The I/O Bottleneck:} Profiling the end-to-end execution time revealed a massive bottleneck in data serialization.
    %    In a test generating 1,533,997 acoustic hits, the internal simulation executed rapidly, but writing the massive dataset to \texttt{ir\_output.json} took over 30 seconds (32.27s single-threaded total vs.
    %    31.28s parallel total).
    %    This demonstrated that standard file I/O operations entirely eclipsed the physics engine's execution time,
    %    necessitating the subsequent implementation of a \texttt{BufWriter} to stage the serialization in RAM.
    %\end{itemize}
\end{document}